% Generated draft: requires mathematical and source review.
\documentclass[11pt]{article}
\usepackage[margin=25mm]{geometry}
\usepackage{fontspec}
\setmainfont{Latin Modern Roman}
\usepackage{amsmath,amssymb,mathtools,unicode-math}
\setmathfont{Latin Modern Math}
\usepackage{enumitem,needspace,xurl,hyperref}
\setlength{\parindent}{0pt}
\setlength{\parskip}{5pt}
\setlength{\emergencystretch}{3em}
\newcommand{\N}{\mathbb{N}}
\newcommand{\Z}{\mathbb{Z}}
\newcommand{\Q}{\mathbb{Q}}
\newcommand{\R}{\mathbb{R}}
\newcommand{\C}{\mathbb{C}}
\newcommand{\F}{\mathbb{F}}
\newcommand{\A}{\mathbb{A}}
\newcommand{\PP}{\mathbb{P}}
\newcommand{\E}{\mathbb{E}}
\newcommand{\eps}{\varepsilon}
\newcommand{\dd}{\,\mathrm{d}}
\begin{document}
\section*{UnsolvedMath ID: 30002787}
\textbf{Source problem code:} catalogue entry
\subsection{Definitions and mathematical statement}
A form is a homogeneous polynomial. The supplied statement does not specify the base field, the common degree, or the convention for scalar coefficients in a Waring decomposition. The following is therefore a conditional formulation whose intended parameters require review.

Fix a field \(K\) and a positive integer \(d\). For a finite set of variables \(X\), write \(K[X]_d\) for the space of homogeneous polynomials of degree \(d\). Under the pure-power convention, define
\[
\operatorname{rk}_{K,d}(G)=\min\left\{s\in\Z_{\geq 0}:G=\sum_{j=1}^{s}L_j^d,\quad L_j\in K[X]_1\right\},
\]
whenever such a decomposition exists. The empty sum gives \(\operatorname{rk}_{K,d}(0)=0\). Another convention allows summands \(c_jL_j^d\), with \(c_j\in K\); the two conventions agree when every scalar in \(K\) has a \(d\)-th root, but the source does not specify which convention is intended.

For the intended field, degree, and rank convention, the additivity question is: for every integer \(r\geq 1\), every collection of finite pairwise disjoint variable sets \(X_1,\ldots,X_r\), and every collection of forms \(F_i\in K[X_i]_d\) for which the ranks are defined, does
\[
\operatorname{rk}_{K,d}\!\left(F_1+\cdots+F_r\right)
=\sum_{i=1}^{r}\operatorname{rk}_{K,d}(F_i)
\]
hold? On the left, the rank is computed in \(K[X_1\cup\cdots\cup X_r]\), so decomposition linear forms may involve variables from several blocks; on the right, each rank is computed in its own variable block.

The common-degree hypothesis is necessary to interpret the sum as a form under this definition, but it is not explicit in the supplied statement. The source also leaves unspecified the allowed fields, characteristics, degrees, and treatment of forms without a Waring decomposition. Consequently, this template does not assert that the intended conjecture ranges over every field or every degree.
\subsection{Short English statement}
The question asks whether the Waring rank of a sum of homogeneous polynomials in disjoint variable sets equals the sum of their individual ranks, even when decompositions of the sum may mix the variable sets. The intended field, common degree, and decomposition convention require clarification. The supplied record labels the problem partially solved; its dated assessment is not independently verified here.
\subsection{Sources}
\begin{itemize}
\item Supplied local UnsolvedMath registry record, permanent ID \(30002787\), title: Additivity of Waring Rank in Independent Variables. Exact source problem code: \texttt{catalogue entry}. The record assigns the year 2015 and supplies no individual proposer.
\item Source report: Mini-Workshop: Ideals of Linear Subspaces, Their Symbolic Powers and Waring Problems (2015), Oberwolfach Reports, pp.\ 489--532, DOI \texttt{10.4171/owr/2015/9}. \url{https://doi.org/10.4171/owr/2015/9}.
\item Supplied literature reference: Carlini, Catalisano, and Chiantini, symmetric Strassen progress. \url{https://doi.org/10.1016/j.jpaa.2014.10.006}.
\item Supplied literature reference: Teitler, sufficient conditions for Strassen additivity. \url{https://doi.org/10.1215/ijm/1488186021}.
\item Supplied literature reference: original OWR report. \url{https://doi.org/10.4171/owr/2015/9}.
\item Literature assessment dated exactly \texttt{2026-08-21}: the supplied curation labels the problem partially solved and reports additivity for several substantial classes, including monomials, binary forms, and linearly computable forms, while describing the general independent-variable conjecture as open. This is supplied dated curation derived from the OWR record, not an independently verified literature assessment; time-sensitive status and the cited claims require independent verification.
\end{itemize}
\end{document}
